\documentclass[a4paper]{article}

\usepackage{graphicx}
\usepackage{amsmath,amssymb}
\usepackage{minted}
\usepackage{multirow}
\usepackage{float}
\usepackage{todonotes}
\usepackage{forest}
\usepackage{xstring}
\usepackage{xcolor}
\usepackage[most]{tcolorbox}
\usepackage{xparse}
\usepackage{natbib}

\usetikzlibrary{graphs,graphdrawing}
\usegdlibrary{layered}

% for external graphviz
\input{/home/donkarlo/Dropbox/repo/nd_language_ai_project/src/nd_language_ai/formal/computer/programming/tex/latex/code_clips/external_graphviz.tex}

% for tags
\input{/home/donkarlo/Dropbox/repo/nd_language_ai_project/src/nd_language_ai/formal/computer/programming/tex/latex/parts/control_sequence/kind/semtag/semtag.tex}

% for links
\input{/home/donkarlo/Dropbox/repo/nd_language_ai_project/src/nd_language_ai/formal/computer/programming/tex/latex/code_clips/seehere.tex}

\title{}
\date{}

\begin{document}
    \maketitle
    
    
    \section{Grammatikthemen im Deutschen}
        
        \begin{itemize}
            \item Verben
            \begin{itemize}
                \item Verbformen und Konjugation
                \begin{itemize}
                    \item Konjugation im Präsens
                    \item starke und schwache Verben
                    \item gemischte Verben
                    \item trennbare und nicht trennbare Verben
                    \item reflexive Verben
                    \item Modalverben
                \end{itemize}
                
                \item Zeiten
                \begin{itemize}
                    \item Präsens
                    \item Perfekt
                    \item Präteritum
                    \item Plusquamperfekt
                    \item Futur I
                \end{itemize}
                
                \item Konjunktiv
                \begin{itemize}
                    \item Konjunktiv II
                    \begin{itemize}
                        \item hätte, wäre, würde
                        \item Konjunktiv II der Vergangenheit
                        \item Konjunktiv II mit Modalverben
                        \item irreale Konditionalsätze
                        \item irreale Wunschsätze
                        \item Höflichkeit
                        \item vorsichtige Aussagen
                    \end{itemize}
                    \item Konjunktiv I
                    \begin{itemize}
                        \item Bildung
                        \item Funktionen
                    \end{itemize}
                \end{itemize}
                
                \item Verbmuster und Ergänzungen
                \begin{itemize}
                    \item Verben mit Akkusativobjekt
                    \item Verben mit Dativobjekt
                    \item Verben mit Dativ- und Akkusativobjekt
                    \item Verben mit Präpositionalobjekt
                    \item Verben mit Pronominaladverbien
                    \item Verben mit festen Ergänzungsmustern
                \end{itemize}
                
                \item Verbale Spezialthemen
                \begin{itemize}
                    \item Bedeutung und Funktion von \textit{werden}
                    \item \textit{brauchen} und \textit{sich lassen}
                    \item Nomen-Verb-Verbindungen
                    \begin{itemize}
                        \item freie Nomen-Verb-Verbindungen
                        \item Funktionsverbgefüge
                        \item figurative Nomen-Verb-Verbindungen
                    \end{itemize}
                \end{itemize}
            \end{itemize}
            
            \item Nomen und Artikel
            \begin{itemize}
                \item Nomen
                \begin{itemize}
                    \item Bedeutung und Form der Nomen
                    \item Genus
                    \item Numerus
                    \item Kasus
                    \item n-Deklination
                    \item Komposita
                \end{itemize}
                
                \item Artikel
                \begin{itemize}
                    \item bestimmte Artikel
                    \item unbestimmte Artikel
                    \item Nullartikel
                    \item Negation mit \textit{kein}
                \end{itemize}
            \end{itemize}
            
            \item Pronomen und Artikelwörter
            \begin{itemize}
                \item Personalpronomen
                \begin{itemize}
                    \item Nominativ
                    \item Akkusativ
                    \item Dativ
                \end{itemize}
                
                \item Possessivartikel
                \item Demonstrativartikel und Demonstrativpronomen
                \item unbestimmte Pronomen und Artikel
                \item Pronominaladverbien
            \end{itemize}
            
            \item Adjektive und Steigerung
            \begin{itemize}
                \item Adjektive
                \item Adjektivdeklination
                \item Komparativ
                \item Superlativ
                \item Vergleichsformen
                \begin{itemize}
                    \item \textit{so \dots wie}
                    \item \textit{als}
                \end{itemize}
            \end{itemize}
            
            \item Kasus
            \begin{itemize}
                \item Nominativ
                \begin{itemize}
                    \item Subjekt
                    \item Prädikativ
                \end{itemize}
                
                \item Akkusativ
                \begin{itemize}
                    \item direktes Objekt
                    \item Akkusativ nach Präpositionen
                    \item temporaler Akkusativ
                \end{itemize}
                
                \item Dativ
                \begin{itemize}
                    \item indirektes Objekt
                    \item Dativ nach Präpositionen
                    \item freier Dativ
                \end{itemize}
                
                \item Genitiv
                \begin{itemize}
                    \item Genitiv als Attribut
                    \item Genitiv bei Namen
                    \item Präpositionen mit Genitiv
                \end{itemize}
            \end{itemize}
            
            \item Präpositionen
            \begin{itemize}
                \item Präpositionen mit Dativ
                \item Präpositionen mit Akkusativ
                \item Präpositionen mit Dativ und Akkusativ
                \item Präpositionen mit Genitiv
                \item lokale Präpositionen
                \item temporale Präpositionen
            \end{itemize}
            
            \item Satzbau
            \begin{itemize}
                \item Wortstellung
                \begin{itemize}
                    \item Aussagesatz
                    \item Fragesatz
                    \item Satzklammer
                    \item Stellung von Satzgliedern
                \end{itemize}
                
                \item Satzarten
                \begin{itemize}
                    \item Aussagesatz
                    \item Fragesatz
                    \item Aufforderungssatz
                    \item Wunschsatz
                    \item Ausrufesatz
                \end{itemize}
                
                \item Satzglieder
                \begin{itemize}
                    \item Prädikat
                    \item Subjekt
                    \item Akkusativobjekt
                    \item Dativobjekt
                    \item Präpositionalobjekt
                    \item Adverbiale
                \end{itemize}
                
                \item Negation
                \begin{itemize}
                    \item \textit{nicht}
                    \item \textit{kein}
                    \item pauschale Negation
                    \item Negationswörter
                \end{itemize}
            \end{itemize}
            
            \item Nebensätze und Satzverbindungen
            \begin{itemize}
                \item Grundtypen
                \begin{itemize}
                    \item dass-Sätze
                    \item ob-Sätze
                    \item Fragesätze als Nebensätze
                    \item Relativsätze
                    \item Infinitivsätze mit \textit{zu}
                \end{itemize}
                
                \item adverbiale Nebensätze
                \begin{itemize}
                    \item kausale Nebensätze
                    \item konzessive Nebensätze
                    \item konditionale Nebensätze
                    \item modale Nebensätze
                    \item temporale Nebensätze
                    \begin{itemize}
                        \item \textit{wenn} und \textit{als}
                        \item \textit{seit[dem]} und \textit{bis}
                        \item \textit{während}
                        \item \textit{nachdem}
                        \item \textit{bevor} und \textit{ehe}
                        \item \textit{sobald}
                        \item \textit{solange}
                    \end{itemize}
                    \item konsekutive Nebensätze
                    \item adversative Nebensätze
                    \item Finalsätze
                    \begin{itemize}
                        \item finale Nebensätze
                        \item finale Infinitivsätze
                        \item \textit{damit}-Sätze
                    \end{itemize}
                \end{itemize}
                
                \item Konjunktionen
                \begin{itemize}
                    \item nebenordnende Konjunktionen
                    \item unterordnende Konjunktionen
                \end{itemize}
            \end{itemize}
            
            \item Passiv
            \begin{itemize}
                \item Vorgangspassiv
                \begin{itemize}
                    \item mit Subjekt
                    \item ohne Subjekt
                    \item Zeiten im Vorgangspassiv
                    \item Vorgangspassiv mit Modalverben
                    \item Vorgangspassiv ohne Täter
                \end{itemize}
                
                \item Zustandspassiv
            \end{itemize}
            
            \item Modalität
            \begin{itemize}
                \item objektive Bedeutungen der Modalverben
                \begin{itemize}
                    \item Möglichkeit
                    \item Fähigkeit
                    \item Erlaubnis
                    \item Verbot
                    \item Notwendigkeit
                    \item Pflicht
                    \item Wunsch
                    \item Wille
                \end{itemize}
                
                \item subjektive Bedeutungen der Modalverben
                \begin{itemize}
                    \item Vermutung
                    \item Unsicherheit
                    \item Gerücht
                    \item Empfehlung
                \end{itemize}
            \end{itemize}
            
            \item Umformungen und Stil
            \begin{itemize}
                \item Nominalisierung
                \item Verbalisierung
                \item Umwandlung adverbialer Präpositionalgruppen in Nebensätze
            \end{itemize}
        \end{itemize}
%    \bibliographystyle{plainnat}
%    \bibliography{/home/donkarlo/Dropbox/repo/nd_shared_research/bibliography/bibliography.bib}
\end{document}